\chapter{Conclusiones y Trabajo Futuro} \label{ch:conclusiones}

\section{Conclusiones técnicas} \label{sec:conclusiones-tecnicas}

El trabajo demuestra la viabilidad de plantear un constructor de gemelos digitales para redes OT/IT a partir de un modelo formal, una interfaz visual y un pipeline de artefactos reproducibles. La arquitectura propuesta permite separar estado visual, modelo técnico, inventario derivado y runtime de laboratorio. Esta separación es fundamental para evitar que el sistema sea solo una herramienta de dibujo o una automatización rígida.

La persistencia granular y el outbox representan una evolución relevante del prototipo. Permiten guardar el trabajo del usuario, desacoplar integraciones externas y procesar tareas derivadas sin bloquear la edición. Además, el uso de Containerlab, Ansible y FRR muestra que el gemelo puede pasar de modelo a laboratorio interactivo.

\section{Conclusiones de ciberseguridad OT} \label{sec:conclusiones-ot}

Desde la perspectiva de ciberseguridad OT, el enfoque aporta valor porque permite analizar sin intervenir directamente sobre producción. La plataforma facilita estructurar inventario, segmentación, comunicaciones, criticidad y evidencias, elementos necesarios para avanzar en ciber-resiliencia y conformidad con marcos como ISA/IEC 62443, NIS2, NIST SP 800-82 y ENS.

La colaboración con AMC Global refuerza la motivación industrial del trabajo. El problema no es abstracto: las organizaciones con varias plantas y procesos críticos necesitan herramientas que permitan comprender y validar sus redes antes de ejecutar cambios que puedan afectar a la continuidad operacional.

\section{Limitaciones del prototipo} \label{sec:limitaciones-prototipo}

El prototipo no sustituye una auditoría formal ni certifica cumplimiento legal. Tampoco incorpora todavía observabilidad completa, análisis Batfish ejecutado de extremo a extremo desde la plataforma, auditoría OPA avanzada, perfiles runtime específicos por fabricante ni RAG normativo con trazabilidad completa de fuentes.

Estas limitaciones delimitan el alcance desarrollado. La aportación principal del TFM no es cerrar todas las capacidades de una plataforma industrial madura, sino diseñar una arquitectura viable, implementar sus bases y validar que el enfoque puede generar inventario, artefactos y evidencias útiles.

\section{Trabajo futuro} \label{sec:trabajo-futuro}

Las líneas de trabajo futuro incluyen:

\begin{itemize}
    \item Integración de observabilidad con LibreNMS y Oxidized para comparar estado diseñado, desplegado y observado.
    \item Auditoría avanzada de cumplimiento mediante OPA y políticas Rego más completas.
    \item Ejecución completa de análisis Batfish desde la plataforma, con consultas de reachability, aislamiento y cambios de configuración.
    \item Incorporación de un sistema RAG normativo con fuentes controladas, citas trazables y separación clara entre evidencia, recomendación y explicación.
    \item Definición de runtime profiles por tipo de activo OT, fabricante, licencia y función.
    \item Gestión multiusuario, autenticación avanzada, gobierno de proyectos y control de permisos por rol.
    \item Validación con escenarios industriales reales o pilotos controlados en colaboración con empresas.
    \item Preparación de una metodología de adopción para organizaciones con inventario incompleto o redes legacy.
\end{itemize}
